\phantomsection
\subsection*{M1C. Sweet Mead}
\addcontentsline{toc}{subsection}{M1C. Sweet Mead}

\textbf{Impressão Geral}: Hidromel com dulçor significativo de mel. Em comparação com um M1B Semi-Sweet Mead, apresenta corpo e caráter de mel mais acentuados. Os níveis de acidez e taninos sustentam o dulçor, mas o dulçor final é um componente importante do sabor. Possui equilíbrio, corpo, final de boca e intensidade de sabor semelhantes aos encontrados em toda a gama de vinhos de sobremesa bem elaborados (como o Ice Wine).

\textbf{Aroma}: O aroma de mel deve ser moderado a intenso, podendo expressar o dulçor e características florais do néctar, refletindo quaisquer variedades de mel declaradas. Buquê de fermentação equilibrado. Hidroméis mais alcoólicos podem apresentar leves notas picantes oriundas do álcool, e os carbonatados tendem a expressar mais aroma.

\textbf{Aparência}: De límpido a brilhante; quanto maior a limpidez, mais desejável. A presença de bolhas, espuma, colarinho ou efervescência deve corresponder ao nível de carbonatação declarado. Cor proveniente das variedades de méis declarados, geralmente de amarelo palha a dourado. Quanto mais vivas e brilhantes as cores, mais desejável.

\textbf{Sabor}: Caráter de mel de moderado a intenso, refletindo as variedades declaradas. Sabores semelhantes aos aromas de mel fermentado em Aroma. Os níveis de dulçor residual variam de médio-alto a alto, com um final de boca notavelmente doce e cheio. A acidez deve ser equilibrada, macia ou vibrante, mas não {crisp}, evitando que o dulçor pareça enjoativo ou dominante. Taninos podem fazer um hidromel suave parecer mais seco. Características ásperas, desagradáveis ou excessivamente sulfurosas ou notas de levedura/autólise são indesejáveis.

\textbf{Sensação na Boca}: Corpo é geralmente de médio-alto a alto. Hidroméis com maior força alcoólica podem apresentar corpo mais alto e maior aquecimento alcoólico. As sensações de corpo não devem vir acompanhadas de dulçor residual enjoativo, mel cru ou não fermentado. A carbonatação segue o nível declarado, podendo ir de tranquilo a espumante.

\textbf{Ingredientes}: Mel de qualquer origem. Espera-se que os méis varietais apresentem as características distintas associadas às variedades declaradas.

\textbf{Comentários}: Sweet Mead é caracterizado por dulçor perceptível. Alguns exemplares podem se assemelhar a vinhos de sobremesa bem elaborados, que possuem acidez e taninos suficientes para equilibrar o dulçor proeminente. A falha mais comum é um dulçor de mel cru, que faz o produto final parecer pouco fermentado.

\textbf{Instruções para Inscrição}: Participantes devem especificar os níveis de dulçor, carbonatação e força alcoólica; neste estilo, o dulçor é de restrito a doce. Variedades de mel podem ser declaradas.

\textbf{Exemplos Comerciais}: Intermiel Benoîte, Moonlight Sensual, Prairie Rose Traditional Sweet, Rabbit's Foot Sweet Mead, Schramm's Mead Sunflowers, White Winter Sweet Mead, Winehaven Stinger.
